\hypertarget{oci-images}{%
\chapter{11. OCI images}\label{oci-images}}

An OCI image is a poor fit for the way Kennel builds a world, and Kennel
runs images anyway. It runs them because images are how workloads are
shipped today, the common denominator nearly every container tool and
deployment assumes, and a confinement framework that could not accept
the format the world has settled on would be confining nothing anyone
runs. Confining them is a concession to adoption, not a claim that the
model is good, and the work is to take the format the world hands Kennel
and confine it as tightly as that format allows.

Every chapter so far assumed the kennel's world was built by Kennel out
of host parts: a fresh tmpfs root, the host's own \texttt{/usr} borrowed
read-only, an execute allowlist naming each binary. The OCI model
changes the ground under all of it. The root is no longer constructed,
it is an image the operator supplied, and the image is untrusted by the
same logic that makes the workload untrusted. What changes when the
floor of the world arrives as a single third-party artefact, and what
has to be rebuilt because it does, is the subject here.

\hypertarget{booting-an-image-without-a-daemon}{%
\section{11.1 Booting an image without a
daemon}\label{booting-an-image-without-a-daemon}}

The obvious way to run an image is to talk to a container daemon, and
Kennel does not, because the daemon is the wrong thing to trust. A
Docker daemon is a root-equivalent orchestration engine, and the two
ways to reach it both fail the premises. Proxying its socket puts a
large stateful protocol parser in the host trusted base in front of a
\texttt{create-container} verb that is host-root. Running a rootless
engine inside the kennel hands the workload \texttt{unshare} and
\texttt{mount}, the two surfaces container escapes most often turn on.
Kennel boots an unpacked rootfs directly under its own namespaces, as
inert filesystem content, and takes neither bridge.

The fetch and the run are split into two verbs so that everything that
parses the image runs where a parser bug is contained:

\begin{kennelcode}
kennel oci build <name> -- <image-ref>   # fetch + unpack into the store
kennel oci run   <name>                  # boot the named entry under policy
\end{kennelcode}

The build pulls and unpacks into a per-operator store; the run resolves
the named entry and boots it under its signed policy. Every parser the
image touches, the registry protocol, the manifest, the tar extraction,
the image's own runtime config, runs inside a confined kennel at
workload authority, so a malicious image breaks a sandbox that was
already disposable and never reaches the daemon
(\texttt{quarantine-the-unsafe}). The trusted base is the size it was.
The split also partitions the policy grammar: an OCI policy carries a
mandatory \texttt{{[}rootfs{]}} block and never a
\texttt{{[}workload{]}}, a normal policy the reverse, and the two are
mutually exclusive at the block level rather than by a runtime check
(\texttt{make-invalid-unrepresentable}). The substrate-replacement
primitive exists only in the grammar the OCI verb consumes, so the
universal model most workloads use never grows it.

\hypertarget{the-image-as-a-declared-substrate}{%
\section{11.2 The image as a declared
substrate}\label{the-image-as-a-declared-substrate}}

Running a third-party image is a trust decision, and Kennel makes the
operator own it rather than making it for them. The operator declares an
image as the execution substrate the way they declare a host directory
writable: an input the policy names and the signature covers. The grant
is loud, in the class of host networking and device passthrough,
carrying a required reason, and the compiler derives a substrate-trust
exposure from its mere presence and surfaces it in the risk report
against that reason (\texttt{footgun-warn-dont-forbid},
\texttt{tell-me-why}). Kennel does not vet the image, launder its
contents, or assert anything about bytes it did not pin (T3.8).

What it does assert is provenance. The build pulls by digest, resolving
a tag to a digest at build time and recording it, so even a tag-built
image freezes to a pinned hash, and a build with no resolvable digest is
refused, because a substrate with no provenance cannot be told apart
from an error (\texttt{refuse-to-start}). The store entry the build
writes is what the run consumes:

\begin{kennelcode}
<store>/<name>/
  rootfs/        the unpacked image filesystem
  config.json    the image's runtime config
  digest         the resolved image@sha256 the build pulled from
  policy.toml    the scaffolded run policy (operator completes + signs)
\end{kennelcode}

The first three are the integrity unit, since the launcher trusts the
config for the entrypoint and env; the policy sits outside it, covered
by its own signature. The signature then covers the substrate, not the
invocation: it pins the image by digest and the policy grants, but not
an exact command line, so the same signed OCI policy can run different
commands where the standard model's signed \texttt{argv} could not.
Provenance is the floor, and content integrity is opt-in hardening
layered above it: a content-addressed store entry the runner verifies
before pivot, and per-file verity that makes every library load and the
config the launcher reads tamper-evident. The operator who needs those
rungs pays for them, the one who does not is not blocked.

\hypertarget{the-overlay}{%
\section{11.3 The overlay}\label{the-overlay}}

An untrusted image cannot simply be the root, because Kennel needs a few
files to win over whatever the image ships and a scratch image ships
almost nothing. The view is built as an overlay. The image is a
read-only lower, never the writable layer, because an image that could
be written would take the workload's writes into the digest-pinned store
entry and dissolve the integrity unit. Above it sits a small Kennel
layer holding the handful of \texttt{/etc} files the resolver and the
persona need, so they win by layer precedence rather than by overwriting
anything, a real file outranking an image symlink at the same path with
no unlink-and-replace. Below it sits a scaffold of empty mountpoint
directories, so a distroless image that ships no \texttt{/proc},
\texttt{/dev}, or \texttt{/tmp} still boots.

\begin{kennelcode}
lowerdir = kennel-etc : image : scaffold   # leftmost wins; the image is a read-only lower
upperdir = <ephemeral | persisted>         # always present; never the image
\end{kennelcode}

Those \texttt{/etc} files are a default and not a boundary:
\texttt{/etc} is writable-through, a workload may copy a file up into
its own layer, and it harms only itself, because what enforces
confinement is the namespace, the BPF, the uid map, Landlock, and
seccomp, never the content of a file.

There is always a writable upper, because an image expects a writable
root, and the operator chooses only whether it survives teardown. The
default discards it. Persisting it is the loud value, because a
persisted upper accumulates divergence outside the integrity ladder, and
the risk report says so. Whole-tree immutability is not the absence of
an upper but a lock expressed over it, which is the next section.

\hypertarget{closure-lock}{%
\section{11.4 Closure-lock}\label{closure-lock}}

The unprivileged build cannot preserve the image's ownership. With no
subuid range it cannot \texttt{chown} extracted files to the image's
uids, so it flattens every inode to the single persona uid, and that
does more than drop root: it erases every ownership boundary inside the
image at once, root from app and app from everything, because there is
now one owner. The workload runs as that owner, so nothing in the tree
is read-only to it by permission, its own binaries included. Left there,
the image under Kennel would be strictly more permissive than the same
image under Docker, where the app uid still cannot write root-owned
\texttt{/usr}.

Closure-lock restores the one boundary the flatten destroyed, that the
executable surface is not self-writable. It cannot be a permission
scheme, because after the flatten there is no foreign owner left for
permissions to protect, and it is not Landlock, because Landlock rights
are additive and a broad write over \texttt{/} cannot be subtracted at
\texttt{/usr}: default-writable-except-the-closure is not a shape a
sub-path rule can express. The lock is a set of read-only mounts over
the merged root, so a write to a locked path fails with \texttt{EROFS}
regardless of owner, and the workload cannot undo it, holding no
\texttt{CAP\_SYS\_ADMIN} in its namespace and barred from \texttt{mount}
by seccomp. The locked set is derived at build and written into the
policy in plain text, where the operator reviews and signs it:

\begin{kennelcode}
[rootfs]
image    = "ghcr.io/org/app@sha256:..."   # pinned by digest
reason   = "vendor image"                  # required; derives T3.8
readonly = ["/usr", "/bin", "/lib"]        # closure-lock: read-only mounts, write => EROFS
writable = ["/usr/lib/python3.12"]         # carve a hole back out; a derived risk
\end{kennelcode}

The build proposes the resurrected boundary and the signature makes it
real.

The derivation turns on one question, whether the image's
\texttt{config.User} is non-root. An image that runs as root declares a
writable substrate by intent, so nothing is locked, which is the same
posture Docker gives a root container and the reason runtime
\texttt{apt} and global \texttt{pip} keep working. That is a deliberate
write-xor-execute waiver: the substrate is granted execute coarsely,
with no per-binary allowlist inside an image, and with nothing locked
the same paths are writable and executable at once, so the separation
the standard run model holds between a read-only \texttt{/usr} and a
default-deny execute does not hold here. A non-root image earns the
separation back: the coarse filesystem closure, \texttt{/usr},
\texttt{/bin}, \texttt{/sbin}, and the library directories the
merged-usr symlinks resolve into, is locked read-only, so its executable
surface stays equal to the pinned image for the session. The estimate is
high-level and has two known holes, an image that drops privilege in its
entrypoint reads as root and gets no lock, and app code outside the
system directories stays writable; the operator closes those by hand
with a writable carve-out for over-reach and an added lock for
under-reach.

\hypertarget{what-holds-and-what-is-waived}{%
\section{11.5 What holds, and what is
waived}\label{what-holds-and-what-is-waived}}

The line worth ending on is which guarantees survive an untrusted root
and which the operator knowingly gives up. Everything the confinement
limb provides survives unchanged, because none of it ever depended on
where the substrate came from: the network namespace and its egress
boundary, the brokered crossings, the proxy and the network policy, the
masked identity, the constructed device allowlist, seccomp, and the
closure-lock boundary for a non-root image. The workload gains no
capability, no \texttt{mount}, no nested user namespace, and the trusted
base does not grow, because the launcher and every image parser run at
workload authority and not in the daemon.

What the operator waives is content. The image supplies its own runtime
closure, the loader and the libraries and the entrypoint, and Kennel
pins the provenance of the whole rather than the identity of each part:
the digest is the assertion, and the dynamic closure beneath it is
unhashed unless the verity rung is enabled. The posture claim is
therefore confinement, not content integrity. An OCI kennel is contained
exactly as tightly as any other, and what the operator accepted, the
moment they named an image as their floor, is that the bytes inside that
floor are the image's word and not Kennel's.
